% TEMPLATE-MILC-v3.tex - plantilla maestra para todas las guias MILC v3
% Ver scripts/generadores/build-guia-milc.py para la lista completa de
% claves esperadas. El script valida que no queden placeholders sin
% reemplazar antes de escribir el archivo final.

\documentclass[11pt,letterpaper]{article}
\usepackage[margin=1.75cm]{geometry}
\usepackage[table]{xcolor}
\usepackage{fontspec}
\usepackage[spanish]{babel}
\usepackage{tikz}
\usepackage[most]{tcolorbox}
\usepackage{tabularx}
\usepackage{array}
\usepackage{fancyhdr}
\usepackage{titlesec}
\usepackage{ragged2e}
\usepackage{enumitem}
\usepackage{microtype}

\setmainfont{Helvetica Neue}
\setsansfont{Helvetica Neue}
\setlength{\parindent}{0pt}
\setlength{\parskip}{6pt}
\hyphenpenalty=200
\exhyphenpenalty=200
\pretolerance=400
\tolerance=2000
\setlength{\emergencystretch}{1em}
\renewcommand{\arraystretch}{1.25}
\setlist{leftmargin=5mm,itemsep=1mm,topsep=1mm}
\newcolumntype{Y}{>{\RaggedRight\arraybackslash}X}

\definecolor{milcMagenta}{HTML}{E6007E}
\definecolor{milcVino}{HTML}{5A0038}
\definecolor{milcTurquesa}{HTML}{0093A5}
\definecolor{milcVerde}{HTML}{58A923}
\definecolor{milcMostaza}{HTML}{E5B400}
\definecolor{milcOcre}{HTML}{B86600}
\definecolor{milcNegro}{HTML}{241816}
\definecolor{milcGris}{HTML}{F7F7F4}
\definecolor{milcLinea}{HTML}{E8E2DD}
\definecolor{milcGradePrimary}{HTML}{0093A5}
\definecolor{milcGradeDark}{HTML}{00525C}
\definecolor{milcGradeSoft}{HTML}{E0F3F5}
\definecolor{milcGradeAccent}{HTML}{CFE7FF}
\definecolor{milcGradeText}{HTML}{FFFFFF}
\color{milcNegro}

\pagestyle{fancy}
\fancyhf{}
\lhead{\small Tecnología e Informática · Grado 8}
\rhead{\small Guía 12 · MILC v3}
\cfoot{\small\thepage}
\renewcommand{\headrulewidth}{0pt}

\newtcolorbox{titlebox}[1]{
  enhanced,colback=#1,colframe=#1,arc=7mm,boxrule=0pt,
  left=7mm,right=7mm,top=3mm,bottom=3mm,
  before skip=8pt,after skip=8pt,
  fontupper=\sffamily\bfseries\Large\color{white}
}

\newtcolorbox{softbox}[3]{
  enhanced,colback=#2,colframe=#1,borderline west={4pt}{0pt}{#1},
  arc=2mm,boxrule=.4pt,left=4mm,right=4mm,top=3mm,bottom=3mm,
  before skip=6pt,after skip=8pt,title={#3},coltitle=white,
  colbacktitle=#1,fonttitle=\sffamily\bfseries,fontupper=\RaggedRight,
  attach boxed title to top left={xshift=3mm,yshift=-2mm},
  boxed title style={arc=2mm, boxrule=0pt}
}

\newtcolorbox{drawbox}[2]{
  enhanced,colback=white,colframe=#1,arc=4mm,boxrule=1pt,
  left=5mm,right=5mm,top=4mm,bottom=4mm,before skip=8pt,after skip=8pt,
  title={#2},coltitle=white,colbacktitle=#1,fonttitle=\sffamily\bfseries,
  fontupper=\RaggedRight,
  attach boxed title to top left={xshift=4mm,yshift=-2mm},
  boxed title style={arc=3mm, boxrule=0pt}
}

\newtcolorbox{infoband}[1]{
  enhanced,colback=#1!8,colframe=#1!50,arc=2mm,boxrule=.5pt,
  left=4mm,right=4mm,top=2mm,bottom=2mm,before skip=4pt,after skip=4pt,
  fontupper=\small\RaggedRight
}

\newcommand{\linead}[1]{\noindent\color{milcLinea}\rule{#1}{0.5pt}\color{milcNegro}}
\newcommand{\respuesta}[1]{\par\vspace{2mm}\foreach \n in {1,...,#1}{\noindent\color{milcLinea}\rule{\linewidth}{0.5pt}\par\vspace{5mm}}\color{milcNegro}}
\newcommand{\checkbox}{\textcolor{milcGradeDark}{\fbox{\phantom{x}}}\hspace{5pt}}

\newcommand{\phasecard}[4]{
\begin{tcolorbox}[enhanced,colback=#3,colframe=#2,arc=3mm,boxrule=.5pt,height=3.05cm,valign=center,halign=center,left=1.5mm,right=1.5mm,top=2mm,bottom=2mm]
{\sffamily\bfseries\fontsize{22}{24}\selectfont\textcolor{#2}{#1}}\par
{\sffamily\bfseries\small #4}
\end{tcolorbox}
}

\begin{document}
\thispagestyle{empty}
\begin{tikzpicture}[remember picture,overlay]
  \fill[white] (current page.south west) rectangle (current page.north east);
  \path[fill=milcGradePrimary,rounded corners=16mm]
    ([xshift=.55cm,yshift=-.55cm]current page.north west) rectangle
    ([xshift=-.55cm,yshift=3.65cm]current page.south east);

  \node[anchor=north west,text=milcGradeText] at ([xshift=2.0cm,yshift=-1.85cm]current page.north west)
    {\sffamily\bfseries\fontsize{14}{16}\selectfont GRADO};
  \node[anchor=north west,text=milcGradeText,opacity=.82] at ([xshift=2.0cm,yshift=-2.75cm]current page.north west)
    {\sffamily\bfseries\fontsize{9}{11}\selectfont SERIE GUÍAS · TECNOLOGÍA E INFORMÁTICA};

  \begin{scope}[shift={([xshift=-3.05cm,yshift=-2.55cm]current page.north east)}]
    \fill[white,opacity=.80] (-.62,-.52) rectangle (.62,.52);
    \draw[milcGradeText,line width=1pt] (-.62,-.52) rectangle (.62,.52);
    \fill[milcGradeDark] (-.42,-.38) rectangle (-.22,.06);
    \fill[milcGradeAccent] (-.10,-.38) rectangle (.10,.24);
    \fill[milcGradePrimary!60!black] (.22,-.38) rectangle (.42,.38);
  \end{scope}

  \node[anchor=west,text=milcGradeText] at ([xshift=1.9cm,yshift=-12.85cm]current page.north west)
    {\sffamily\bfseries\fontsize{124}{124}\selectfont 8};
  \draw[milcGradeText,line width=1.7pt]
    ([xshift=4.52cm,yshift=-11.68cm]current page.north west) circle (.13cm);

  \node[anchor=west,text=milcGradeText,text width=8.7cm,align=left] at ([xshift=10.35cm,yshift=-11.6cm]current page.north west)
    {
     {\sffamily\bfseries\fontsize{19}{22}\selectfont Tablas de verdad y condiciones compuestas: AND, OR y NOT en decisiones}\\[4mm]
     {\sffamily\bfseries\fontsize{23}{26}\selectfont Guía 12-8-TIC}\\[2mm]
     {\sffamily\fontsize{13}{16}\selectfont Periodo Segundo}\\[5mm]
     {\sffamily\bfseries\fontsize{11}{13}\selectfont Institución Educativa Sor María Juliana}\\[1mm]
     {\sffamily\fontsize{10}{12}\selectfont Autor: Dr. Álvaro Cárdenas Orozco}
    };

  \path[fill=milcGradeSoft,rounded corners=7mm]
    ([xshift=1.35cm,yshift=.85cm]current page.south west) rectangle
    ([xshift=-1.35cm,yshift=3.15cm]current page.south east);
  \draw[milcGradeDark,line width=.65pt,rounded corners=7mm]
    ([xshift=1.35cm,yshift=.85cm]current page.south west) rectangle
    ([xshift=-1.35cm,yshift=3.15cm]current page.south east);
  \node[anchor=center,text=milcNegro,text width=17.0cm,align=center] at ([yshift=2.0cm]current page.south)
    {\sffamily\fontsize{10}{12}\selectfont Metodología MILC · Apertura ancestral · Pensamiento computacional · Triángulo Dussel-Estoicismo-Floridi\\
    \textcolor{milcGradeDark}{\bfseries Producto:} 3 tablas de verdad de expresiones compuestas + 1 explicación de un caso real donde una ``Y'' se confundió con una ``O''.};
\end{tikzpicture}
\mbox{}
\newpage

\section*{Apertura: Tablas de verdad y condiciones compuestas: AND, OR y NOT en decisiones}

\begin{softbox}{milcGradeDark}{milcGradeSoft}{Saber ancestral}
El \textbf{juego de ``verdadero o falso''} que las abuelas usaban para enseñar a no mentir: \textit{``si dices A y dices B, ¿ambas pueden ser verdad al tiempo? ¿o se contradicen?''}. La tabla de verdad sistematiza ese juego ancestral. Antes de existir la lógica formal, las abuelas ya entrenaban a los niños en consistencia entre afirmaciones.
\end{softbox}

\begin{softbox}{milcTurquesa}{milcGris}{Saber contemporáneo}
Una \textbf{tabla de verdad} enumera todos los casos posibles de una expresión booleana y muestra el resultado de cada uno. Para n variables, hay 2\^{}n casos. Sirve para verificar que la lógica es correcta antes de programar.
\end{softbox}

\begin{softbox}{milcMostaza}{milcGris}{Pregunta-puente}
\textit{¿Qué del juego de las abuelas —el cuidado de la consistencia entre afirmaciones— podemos llevar a las tablas de verdad de la programación?}
\end{softbox}

\begin{softbox}{milcVerde}{milcGris}{Saber y saber hacer}
Al terminar podrás \textbf{construir tablas de verdad} de hasta 3 variables, \textbf{evaluar expresiones compuestas} caso por caso, \textbf{detectar equivalencias lógicas} (cuando dos expresiones distintas dan los mismos resultados) y \textbf{simplificar} expresiones complejas.
\end{softbox}

\small
\begin{tabularx}{\textwidth}{>{\bfseries}p{3.0cm}Y}
\rowcolor{milcGradePrimary!12}Autor & Dr. Álvaro Cárdenas Orozco\\
DBA & Construye tablas de verdad para verificar consistencia lógica de expresiones compuestas con AND, OR y NOT.\\
Referentes & MEN Tecnología · ICFES Pensamiento Lógico · Modelo MILC · Floridi (impacto del bien/mal informacional) · Epicteto (acción coherente)\\
Producto final & 3 tablas de verdad de expresiones compuestas + 1 explicación de un caso real donde una ``Y'' se confundió con una ``O''.\\
\end{tabularx}
\normalsize

\begin{titlebox}{milcGradeDark}
Ruta MILC: así vamos a aprender
\end{titlebox}
\begin{tabularx}{\textwidth}{>{\centering\arraybackslash}X>{\centering\arraybackslash}X>{\centering\arraybackslash}X>{\centering\arraybackslash}X}
\phasecard{1}{milcVerde}{milcGris}{Escuta\\[-1mm]\small Observo, escucho y pregunto.} &
\phasecard{2}{milcTurquesa}{milcGris}{Sistematización\\[-1mm]\small Pensamiento computacional.} &
\phasecard{3}{milcMagenta}{milcGris}{Praxis\\[-1mm]\small Construyo y aplico.} &
\phasecard{4}{milcVino}{milcGris}{Evaluación\\[-1mm]\small Reflexión liberadora.}
\end{tabularx}


\begin{titlebox}{milcVerde}
Fase 1 · Escuta
\end{titlebox}

\begin{softbox}{milcVerde}{milcGris}{Escena para escuchar}
Diana hace la regla con su mamá: \textit{``Si gano el partido Y apruebo el examen, salgo a celebrar''}. El sábado gana el partido pero no aprueba. Diana asume que igual sale (porque cumplió una). La mamá dice no — la regla era \textit{``Y''}, no \textit{``O''}. Diana descubre que confundió un AND con un OR. ¿Cómo se evita ese error? Con tabla de verdad.
\end{softbox}

\checkbox Recuerdo una situación familiar donde se confundió ``Y'' con ``O'' (o viceversa).\\
\checkbox Identifico la diferencia entre conjunción (AND) y disyunción (OR) en frases reales.\\
\checkbox Distingo cuándo una tabla de verdad ayuda a aclarar una regla.

\textbf{Una situación donde se confundió Y con O en mi entorno fue:}\\[1mm]
\linead{\linewidth}\\[1mm]
\linead{\linewidth}

\textbf{Si hubiéramos hecho la tabla de verdad, habríamos visto que:}\\[1mm]
\linead{\linewidth}\\[1mm]
\linead{\linewidth}

\textbf{Una regla familiar/escolar que mejoraría con tabla de verdad es:}\\[1mm]
\linead{\linewidth}\\[1mm]
\linead{\linewidth}

\begin{infoband}{milcVerde}
\textbf{Saber más.} La tabla de verdad de \texttt{A AND B} solo tiene UN caso verdadero (cuando ambas son V). La de \texttt{A OR B} tiene tres casos verdaderos (todas menos cuando ambas son F). Por eso la diferencia importa: AND es exigente, OR es permisivo.
\end{infoband}

\begin{titlebox}{milcTurquesa}
Fase 2 · Sistematización con pensamiento computacional
\end{titlebox}

\begin{softbox}{milcTurquesa}{milcGris}{Cuatro pilares aplicados al tema}
El pensamiento computacional descompone la expresión en variables, reconoce patrones de evaluación, abstrae la regla \textit{enumerar todos los casos} y diseña la tabla como algoritmo verificador.
\end{softbox}

\small
\begin{tabularx}{\textwidth}{>{\bfseries\RaggedRight\arraybackslash}p{3.6cm}Y}
\rowcolor{milcTurquesa!90}\color{white}Pilar & \color{white}\bfseries Aplicación al tema\\
1. Descomposición & Descompón: ¿cuántas variables? Si son n, hay 2\^{}n casos. Listar todas las combinaciones.\\
2. Reconocimiento de patrones & Patrones: AND es V solo si ambas V. OR es V si al menos una V. NOT invierte cada caso.\\
3. Abstracción & Abstracción: la tabla es la función booleana en forma de inventario completo.\\
4. Algoritmo conceptual & Algoritmo: \textbf{1)} listar variables · \textbf{2)} crear todas las combinaciones (2\^{}n filas) · \textbf{3)} evaluar la expresión en cada fila · \textbf{4)} marcar V o F · \textbf{5)} contar cuántos V (revela el carácter de la expresión).\\
\end{tabularx}
\normalsize

\begin{drawbox}{milcTurquesa}{Anatomía de una tabla de verdad de 3 variables}
\textbf{Variables:} A, B, C (3 variables). \\
\textbf{Filas:} 2\^{}3 = 8 combinaciones posibles. \\
\textbf{Orden estándar:} de FFF a VVV (binario ascendente). \\
\textbf{Columnas auxiliares:} resultados parciales de subexpresiones. \\
\textbf{Columna final:} resultado de la expresión completa. \\
\textbf{Verificación:} cada fila independiente; todas se evalúan.
\end{drawbox}

\begin{softbox}{milcMostaza}{milcGris}{Errores comunes y su corrección}
\begin{itemize}
\item Olvidar combinaciones $\rightarrow$ siempre 2\^{}n filas; ni más ni menos.
\item Evaluar mal el orden de operadores $\rightarrow$ NOT primero, luego AND, luego OR (igual que en código).
\item Confundir equivalencia con igualdad $\rightarrow$ \texttt{A AND B} y \texttt{B AND A} son equivalentes (mismas Vs); \texttt{A AND B} y \texttt{A OR B} no lo son.
\item No usar columnas auxiliares en expresiones largas $\rightarrow$ ayuda a no perderse.
\item Saltarse casos por flojera $\rightarrow$ la fila que omites es justo donde puede estar el error.
\end{itemize}
\end{softbox}

\textbf{La tabla de verdad de \texttt{A AND B} en mis palabras significa:}\\[1mm]
\linead{\linewidth}\\[1mm]
\linead{\linewidth}

\textbf{Una expresión equivalente a \texttt{NOT (A AND B)} sería:}\\[1mm]
\linead{\linewidth}\\[1mm]
\linead{\linewidth}

\begin{titlebox}{milcMagenta}
Fase 3 · Praxis con pensamiento computacional
\end{titlebox}

\begin{softbox}{milcMagenta}{milcGris}{Construyendo el producto con los cuatro pilares}
Construir tablas de verdad es ejercicio de phronesis lógica. Decomponemos la expresión, reconocemos patrones, abstraemos la regla y seguimos un algoritmo verificable.
\end{softbox}

\small
\begin{tabularx}{\textwidth}{>{\bfseries\RaggedRight\arraybackslash}p{3.6cm}Y}
\rowcolor{milcMagenta!90}\color{white}Pilar & \color{white}\bfseries Aplicación al producto\\
Descomposición & Tabla: 1) variables · 2) combinaciones · 3) subexpresiones · 4) resultado final · 5) interpretación.\\
Patrones & Patrones de tablas profesionales: orden binario ascendente, columnas auxiliares, resultado final destacado.\\
Abstracción & Función esencial: la tabla CERTIFICA que la expresión funciona en todos los casos posibles.\\
Algoritmo de armado & Pasos: identificar variables $\rightarrow$ enumerar combinaciones $\rightarrow$ evaluar paso a paso $\rightarrow$ contar Vs $\rightarrow$ interpretar carácter.\\
\end{tabularx}
\normalsize

\begin{drawbox}{milcMagenta}{Checklist de las 3 tablas de verdad}
\checkbox Tabla 1: expresión \texttt{A AND (B OR C)}.\\
\checkbox Tabla 2: expresión \texttt{(A OR B) AND NOT C}.\\
\checkbox Tabla 3: expresión a tu elección con 3 variables.\\
\checkbox 8 filas en cada (2\^{}3 casos).\\
\checkbox Columnas auxiliares para subexpresiones.\\
\checkbox Caso real explicado donde se confundió Y con O.
\end{drawbox}

\begin{softbox}{milcMagenta}{milcGris}{Plantilla de trabajo}
\textbf{Expresión:} \texttt{(A AND B) OR NOT C} (ejemplo). \\
\textbf{Variables:} A, B, C. \\
\textbf{Filas:} 8 combinaciones de FFF a VVV. \\
\textbf{Columnas:} A, B, C, A AND B, NOT C, (A AND B) OR NOT C. \\
\textbf{Evaluación:} fila por fila, sin saltarse. \\
\textbf{Caso real:} \textit{``En mi familia, la regla era X; alguien la entendió como Y. La tabla muestra que la regla correcta era \_\_\_.''}
\end{softbox}

\textbf{Mis 3 tablas de verdad + 1 caso real:}\\[1mm]
\linead{\linewidth}\\[1mm]
\linead{\linewidth}\\[1mm]
\linead{\linewidth}

\begin{titlebox}{milcMostaza}
Producto de la sesión
\end{titlebox}
\textbf{3 tablas de verdad + 1 caso real explicado}.
\begin{softbox}{milcMostaza}{milcGris}{Criterios de calidad}
\begin{itemize}
\item 3 tablas con 3 variables (8 filas cada una).
\item Columnas auxiliares para subexpresiones.
\item Resultado final marcado claramente.
\item 1 caso real bien explicado.
\item Reflexión sobre cómo la tabla aclaró la regla.
\end{itemize}
\end{softbox}

\begin{titlebox}{milcVino}
Fase 4 · Evaluación liberadora — cinco dimensiones
\end{titlebox}

\begin{softbox}{milcVino}{milcGris}{Sentido de la evaluación}
Evaluar no es solo revisar una nota. Es reconocer qué aprendí, qué cambió en mí, qué emociones manejé, cómo me relaciono con la ciudad y la comunidad, y qué le dejaré a quien viene detrás.
\end{softbox}

\textbf{Mi reflexión liberadora en cinco dimensiones.} Completa cada frase iniciada:

\small
\begin{tabularx}{\textwidth}{>{\bfseries\RaggedRight\arraybackslash}p{3.4cm}Y>{\RaggedRight\arraybackslash}p{4.6cm}}
\rowcolor{milcVino}\color{white}Dimensión & \color{white}\bfseries Mi respuesta (frame starter) & \color{white}\bfseries Algo que descubrí\\
1. Personal & Lo que más cambió en mí fue \linead{6.5cm} & \linead{4.2cm}\\
2. Emocional & La emoción más fuerte fue \linead{2.5cm} y la regulé \linead{3cm} & \linead{4.2cm}\\
3. Ciudadana & En democracia esto importa porque \linead{6cm} & \linead{4.2cm}\\
4. Local (Cartago) & En mi barrio sería útil para \linead{6.5cm} & \linead{4.2cm}\\
5. Intergeneracional & A mi abuela/abuelo le diría \linead{6.5cm} & \linead{4.2cm}\\
\end{tabularx}
\normalsize

\begin{infoband}{milcVino}
\textbf{Para la próxima sustentación.} Vuelve a esta tabla en seis meses. Verás que las respuestas cambian — no porque lo aprendido cambie, sino porque tú cambias. La evaluación liberadora es una conversación contigo a través del tiempo.
\end{infoband}


\begin{titlebox}{milcGradeDark}
Triángulo de pensamiento · cierre filosófico
\end{titlebox}

\begin{softbox}{milcVino}{milcGris}{Dussel · lente del nosotros}
\textit{``El Otro es el origen absoluto de la responsabilidad ética.''}

Una regla mal aplicada (Y confundido con O) puede excluir injustamente o incluir indebidamente al Otro. Diana fue excluida del paseo por una regla que ella entendía distinto. La tabla de verdad es ética del lenguaje claro.

\textbf{¿Qué regla en tu entorno se interpreta distinto según la persona? ¿Cómo se aclararía con tabla de verdad?}
\end{softbox}
\linead{\linewidth}\\[1mm]
\linead{\linewidth}

\begin{softbox}{milcMostaza}{milcGris}{Estoicismo · Epicteto · cuidado interior}
\textit{``No es lo que dices, es lo que haces lo que define quién eres.''}

Si afirmas A y luego afirmas \textit{NOT A}, te contradices. La tabla de verdad detecta contradicciones imposibles. La coherencia entre lo que dices y lo que haces se prueba en la lógica del día a día.

\textbf{¿Hay alguna afirmación tuya reciente que con tabla de verdad se vea contradictoria?}
\end{softbox}
\linead{\linewidth}\\[1mm]
\linead{\linewidth}

\begin{softbox}{milcTurquesa}{milcGris}{Floridi · lente de la infoesfera}
\textit{``El bien y el mal informacional se miden por su impacto en seres y datos.''}

Una expresión lógica mal escrita en software de pruebas médicas puede dar diagnósticos falsos. El bien informacional empieza con tablas de verdad correctas. La precisión lógica es ética antes de ser técnica.

\textbf{Si tu lógica se aplicara a un caso real (alarma, filtro de spam, regla escolar), ¿pasaría tu tabla de verdad la auditoría?}
\end{softbox}
\linead{\linewidth}\\[1mm]
\linead{\linewidth}

\begin{titlebox}{milcVino}
Compromiso final
\end{titlebox}

\small
\begin{tabularx}{\textwidth}{>{\RaggedRight\arraybackslash}p{3.0cm}YYY}
\rowcolor{milcVino}\color{white}\bfseries Criterio & \color{white}\bfseries Lo logré & \color{white}\bfseries En proceso & \color{white}\bfseries Necesito apoyo\\
Comprensión del tema & Explico ideas centrales con mis palabras. & Reconozco ideas pero falta precisión. & Necesito repasar conceptos base.\\
Pensamiento computacional & Apliqué los 4 pilares al diseñar mi producto. & Apliqué algunos pilares. & Necesito releer la fase 2.\\
Aplicación y producto & Mi producto cumple los criterios y se puede revisar. & Producto funcional con ajustes pendientes. & Aún en construcción.\\
Reflexión 5 dimensiones & Respondí cada dimensión con honestidad. & Respondí algunas con detalle. & Mis respuestas son muy generales.\\
Triángulo de pensamiento & Conecté las 3 voces con mi trabajo real. & Conecté al menos una. & No relaciono las voces con mi proceso.\\
\end{tabularx}
\normalsize

\textbf{Hoy entendí que:}\\[1mm]
\linead{\linewidth}\\[1mm]
\linead{\linewidth}

\textbf{Mi compromiso para la próxima sesión es:}\\[1mm]
\linead{\linewidth}\\[1mm]
\linead{\linewidth}

\begin{softbox}{milcMostaza}{milcGris}{Cita de despedida · Marco Aurelio}
\textit{``El obstáculo en el camino se vuelve el camino. Lo que estorba al camino se vuelve camino.''}\\[1mm]
Cada error de hoy, bien observado, se convierte en pieza del camino que pisarás mañana.
\end{softbox}

\small
\begin{tabularx}{\textwidth}{>{\bfseries}p{3.6cm}Y}
\rowcolor{milcGradeDark!12}Nombre completo: & \linead{10cm}\\
Fecha: & \linead{4cm} \quad Curso/grupo: \linead{4cm}\\
Mi firma: & \linead{9cm}\\
\end{tabularx}
\normalsize

\begin{infoband}{milcGradeDark}
\textbf{Próxima sesión.} Aprenderemos pseudocódigo y diagramas de flujo. Las tablas de verdad se vuelven mapas visuales de la decisión.
\end{infoband}

\end{document}
